\documentclass[12pt,a4paper]{article}

% Encoding and fonts — XeLaTeX/LuaLaTeX
\usepackage{fontspec}
\newfontfamily\hebrewfont[Script=Hebrew]{Noto Serif Hebrew}
\newcommand{\heb}[1]{{\hebrewfont #1}}
\newfontfamily\igfont[Ligatures=TeX]{Noto Serif}
\newcommand{\igtext}[1]{{\igfont #1}}

% Page layout
\usepackage[top=1in, bottom=1in, left=1in, right=1in]{geometry}

% Typography
\usepackage{microtype}
\usepackage{parskip}

% Language
\usepackage[english]{babel}

% Headers and footers
\usepackage{fancyhdr}
\pagestyle{fancy}
\fancyhf{}
\fancyhead[L]{\leftmark}
\fancyhead[R]{\thepage}
\fancyfoot[C]{\thepage}
\renewcommand{\headrulewidth}{0.4pt}
\renewcommand{\footrulewidth}{0pt}

% Section styling
\usepackage{titlesec}
\titleformat{\section}{\Large\bfseries}{\thesection}{1em}{}
\titleformat{\subsection}{\large\bfseries}{\thesubsection}{1em}{}
\titleformat{\subsubsection}{\normalsize\bfseries}{\thesubsubsection}{1em}{}
\titlespacing*{\section}{0pt}{2.5ex plus 1ex minus .2ex}{1.5ex plus .2ex}
\titlespacing*{\subsection}{0pt}{2ex plus 1ex minus .2ex}{1ex plus .2ex}
\titlespacing*{\subsubsection}{0pt}{1.5ex plus 1ex minus .2ex}{0.8ex plus .2ex}

% List formatting
\usepackage{enumitem}
\setlist[itemize]{topsep=0.5ex, itemsep=0.25ex, parsep=0pt}
\setlist[enumerate]{topsep=0.5ex, itemsep=0.25ex, parsep=0pt}

% Essential packages
\usepackage{graphicx}
\usepackage[table]{xcolor}
\usepackage{booktabs}
\usepackage{array}
\usepackage{tabularx}
\usepackage{longtable}
\usepackage{amsmath}
\usepackage{amssymb}
\usepackage[normalem]{ulem}
\rowcolors{2}{gray!10}{white}
\renewcommand{\arraystretch}{1.3}

% Cross-references
\usepackage{hyperref}
\usepackage{cleveref}

% Custom commands
\newcommand{\pilcrow}{\S}

% Hyperref setup
\hypersetup{
    colorlinks=true,
    linkcolor=blue,
    filecolor=magenta,
    urlcolor=cyan,
}

% Code listings
\usepackage{listings}
\definecolor{codebg}{RGB}{248,248,248}
\definecolor{codeframe}{RGB}{200,200,200}
\definecolor{codekeyword}{RGB}{0,0,180}
\definecolor{codecomment}{RGB}{0,128,0}
\definecolor{codestring}{RGB}{163,21,21}
\lstset{
    basicstyle=\ttfamily\small,
    breaklines=true,
    breakatwhitespace=false,
    frame=single,
    framerule=0.5pt,
    rulecolor=\color{codeframe},
    backgroundcolor=\color{codebg},
    keepspaces=true,
    numbers=left,
    numberstyle=\tiny\color{gray},
    numbersep=8pt,
    showstringspaces=false,
    tabsize=4,
    xleftmargin=1em,
    xrightmargin=1em,
    aboveskip=1em,
    belowskip=1em,
    keywordstyle=\color{codekeyword}\bfseries,
    commentstyle=\color{codecomment}\itshape,
    stringstyle=\color{codestring},
    literate={_}{{\textunderscore}}1
             {-}{{\textminus}}1
             {<}{{\textless}}1
             {>}{{\textgreater}}1
             {~}{{\textasciitilde}}1
             {^}{{\textasciicircum}}1
}

% YAML language definition for listings
\lstdefinelanguage{YAML}{
    keywords={true,false,null,y,n},
    keywordstyle=\color{blue}\bfseries,
    sensitive=false,
    comment=[l]{\#},
    morecomment=[s]{/*}{*/},
    commentstyle=\color{gray}\ttfamily,
    stringstyle=\color{red}\ttfamily,
    morestring=[b]'',
    morestring=[b]'',
}

% TOML language definition for listings
\lstdefinelanguage{TOML}{
    keywords={true,false},
    keywordstyle=\color{blue}\bfseries,
    sensitive=true,
    comment=[l]{\#},
    commentstyle=\color{gray}\ttfamily,
    stringstyle=\color{red}\ttfamily,
    morestring=[b]'',
    morestring=[b]'',
}

% Dockerfile language definition for listings
\lstdefinelanguage{Dockerfile}{
    keywords={FROM,RUN,CMD,LABEL,EXPOSE,ENV,ADD,COPY,ENTRYPOINT,VOLUME,USER,WORKDIR,ARG,ONBUILD,STOPSIGNAL,HEALTHCHECK,SHELL},
    keywordstyle=\color{blue}\bfseries,
    sensitive=false,
    comment=[l]{\#},
    commentstyle=\color{gray}\ttfamily,
    stringstyle=\color{red}\ttfamily,
    morestring=[b]'',
    morestring=[b]'',
}

% JSON language definition for listings
\lstdefinelanguage{JSON}{
    keywords={true,false,null},
    keywordstyle=\color{blue}\bfseries,
    sensitive=true,
    commentstyle=\color{gray}\ttfamily,
    stringstyle=\color{red}\ttfamily,
    morestring=[b]'',
}

% Code listing float environment
\usepackage{float}
\newfloat{listing}{htbp}{lol}[section]
\floatname{listing}{Listing}

% Admonition boxes
\usepackage{tcolorbox}
\tcbuselibrary{skins,breakable}

% Admonition style definitions
\newtcolorbox{admonitionbox}[2][]{
    enhanced,
    breakable,
    colback=#2!5,
    colframe=#2!80!black,
    fonttitle=\bfseries,
    title=#1,
    left=2mm,
    right=2mm,
}

% Imscription tuple boxes
\newtcolorbox{imscriptionbox}{
    enhanced,
    colback=blue!5,
    colframe=blue!75!black,
    fontupper=\large,
    center upper,
    boxrule=1pt,
    arc=4pt,
}

\begin{document}

\title{The Aether and Its Vessel: A Formal Proof that  Finds Its Perfect Vessel in}
\date{\today}

\maketitle

\textbf{A Structural Demonstration Without Recourse to Imscribing Primitives}

\begin{center}
\rule{0.5\textwidth}{0.4pt}
\end{center}

\subsection{Abstract}

We present a formal proof that the exceptional Lie group $E_8$ --- the maximal exceptional structure, which we term the \emph{Aether} --- necessarily finds its completion in $G_2$, the smallest exceptional Lie group, which we term the \emph{Vessel}. The proof proceeds entirely within the established framework of Lie theory, root system analysis, and octonion algebra. We demonstrate three core results: (1) $G_2 \subset E_8$ as a maximal subgroup, establishing that $E_8$ \emph{contains} $G_2$; (2) $G_2$ is the unique minimal exceptional Lie group --- the irreducible structural seed from which all exceptional phenomena emerge; (3) the containment is \emph{perfect} in the sense that the root system of $E_8$ decomposes in such a way that $G_2$ appears as the ineliminable core, the minimal structure that can receive and hold the full octonionic non-associativity that $E_8$ requires. The proof reveals $G_2$ as the \emph{perfect vessel}: the smallest structure capable of holding all the exceptionality of $E_8$, and $E_8$ as the \emph{Aether}: the maximal unfolding of what $G_2$ already contains in seed form.

\begin{center}
\rule{0.5\textwidth}{0.4pt}
\end{center}

\subsection{Introduction: The Exceptional Chain}

The exceptional Lie groups form a chain of increasing dimension and complexity:

\[G_2 \subset F_4 \subset E_6 \subset E_7 \subset E_8\]

This chain exhibits a remarkable property: each containment is a maximal subgroup embedding, and the entire chain is generated by successive extensions of the octonion algebra $\mathbb{O}$. The smallest member, $G_2$, has dimension 14 and rank 2. The largest, $E_8$, has dimension 248 and rank 8. The ratio of their dimensions is $248/14 \approx 17.7$, yet structurally they are separated by a single fundamental step: the passage from a bounded, finite-dimensional crystalline structure to an unbounded, infinite-dimensional field of possibility.

What we establish here is that this chain is not merely a sequence of independent structures but a single reality viewed at different scales: $G_2$ is the \emph{vessel} --- the minimal closure capable of receiving octonionic non-associativity --- and $E_8$ is the \emph{aether} --- the maximal unfolding of that same non-associativity into a self-consistent gauge structure. The proof demonstrates that the relationship is necessary, not contingent.

\begin{center}
\rule{0.5\textwidth}{0.4pt}
\end{center}

\subsection{Octonions and the Automorphism Group $G_2$}

\subsubsection{The Octonion Algebra}

The octonions $\mathbb{O}$ constitute the largest normed division algebra, of real dimension 8. Unlike the quaternions $\mathbb{H}$, the octonions are non-associative. Their multiplication is governed by the Cayley-Dickson construction applied to $\mathbb{H}$, yielding a non-associative alternative algebra with a multiplicative norm.

The non-associativity of $\mathbb{O}$ is not a defect but the essential feature that makes exceptional structures possible. If the octonions were associative, the exceptional Lie algebras would collapse to classical series. The non-associativity is the \emph{generative tension} from which all exceptional phenomena arise.

\subsubsection{$G_2$ as $\text{Aut}(\mathbb{O})$}

$G_2$ is defined as the automorphism group of the octonions:

\[G_2 = \{ \phi \in \text{GL}(8,\mathbb{R}) \mid \phi(xy) = \phi(x)\phi(y) \text{ for all } x,y \in \mathbb{O} \}\]

This is the group of all linear transformations of $\mathbb{O}$ that preserve the octonion multiplication. $G_2$ is the unique compact, connected, simply-connected, simple Lie group of dimension 14 and rank 2.

The key structural fact is this: $G_2$ preserves precisely the non-associativity of $\mathbb{O}$. It does not eliminate it; it stabilizes it. The 14 dimensions of $G_2$ correspond to the 14 independent ways in which the octonion product can be rotated while preserving its non-associative character. $G_2$ is thus the \emph{minimal enclosure} of non-associativity --- the smallest symmetry group that can hold the octonion cross product as an invariant.

\subsubsection{The Seven-Dimensional Representation}

$G_2$ acts irreducibly on the 7-dimensional space of pure imaginary octonions $\text{Im}(\mathbb{O})$. This 7-dimensional representation is the smallest non-trivial representation of $G_2$ and is intimately connected to the geometry of the 7-sphere $S^7$, which is the unit sphere in $\mathbb{O}$. The action of $G_2$ on $S^7$ preserves a non-degenerate 3-form $\varphi$ and its Hodge dual 4-form $\star\varphi$, which define a $G_2$-structure --- a reduction of the structure group of the tangent bundle from $\text{SO}(7)$ to $G_2$.

This geometry is the template for all subsequent exceptional structures: it is the \emph{vessel pattern} that, when extended, generates $F_4$, $E_6$, $E_7$, and ultimately $E_8$.

\begin{center}
\rule{0.5\textwidth}{0.4pt}
\end{center}

\subsection{$E_8$ as the Maximal Unfolding}

\subsubsection{The Root System of $E_8$}

The $E_8$ root system is the largest exceptional root system, consisting of 240 roots in 8-dimensional Euclidean space. These roots are the vertices of the $4_{21}$ polytope, a remarkable 8-dimensional semi-regular polytope with 240 vertices, 6,720 edges, and rich internal structure.

The 240 roots of $E_8$ can be constructed as follows. Let $\{e_i\}_{i=1}^8$ be an orthonormal basis of $\mathbb{R}^8$. The root system $\Phi_{E_8}$ consists of:

\begin{enumerate}
    \item All vectors of the form $\pm e_i \pm e_j$ for $1 \leq i < j \leq 8$ (112 roots);
    \item All vectors of the form $\frac{1}{2}(\pm e_1 \pm e_2 \pm \cdots \pm e_8)$ with an even number of minus signs (128 roots).
\end{enumerate}

Together these yield $112 + 128 = 240$ roots. The Coxeter number of $E_8$ is 30, and the order of the Weyl group is $696,729,600 = 2^{14} \cdot 3^5 \cdot 5^2 \cdot 7$.

\subsubsection{Maximal Subalgebras and the $G_2$ Embedding}

Within the $E_8$ Lie algebra, there exist maximal subalgebras of exceptional type. The maximal subalgebras of $E_8$ have been classified by Dynkin (1957). Among them, one finds:

\[G_2 \times F_4 \subset E_8\]

as a maximal subalgebra. This is the key embedding: $G_2$ appears as a maximal subgroup within $E_8$, paired with $F_4$ in a product structure. More fundamentally, one can trace through the chain:

\[G_2 \subset F_4 \subset E_6 \subset E_7 \subset E_8\]

where each step is itself a maximal embedding. In particular, $G_2 \subset F_4$ is a maximal embedding (via the ''''quaternionic'''' magic square construction), and $F_4 \subset E_6 \subset E_7 \subset E_8$ follows via the standard chain of exceptional groups.

The existence of this chain establishes that $G_2$ is not merely \emph{contained in} $E_8$ as an arbitrary subgroup, but is embedded as an \emph{irreducible structural core} --- the minimal exceptional seed from which the entire tree grows.

\subsubsection{The 248-Dimensional Adjoint Representation}

The adjoint representation of $E_8$, of dimension 248, decomposes under the $G_2 \times F_4$ subalgebra. The decomposition reveals how $E_8$ ''''contains'''' $G_2$:

Under $G_2 \times F_4 \subset E_8$, the adjoint representation decomposes as:

\[\mathbf{248} \to (\mathbf{14}, \mathbf{1}) \oplus (\mathbf{1}, \mathbf{52}) \oplus (\mathbf{7}, \mathbf{26})\]

where $\mathbf{14}$ is the adjoint of $G_2$, $\mathbf{52}$ is the adjoint of $F_4$, $\mathbf{7}$ is the fundamental representation of $G_2$, and $\mathbf{26}$ is the fundamental representation of $F_4$.

This decomposition is remarkable: the $G_2$ adjoint appears as a direct summand, untouched by $F_4$. The $G_2$ structure is \emph{preserved intact} within $E_8$. It is not modified, augmented, or approximated --- it is present in its exact 14-dimensional form. The $G_2$ that lives within $E_8$ is precisely the same $G_2$ that stands alone as the automorphism group of the octonions.

\begin{center}
\rule{0.5\textwidth}{0.4pt}
\end{center}

\subsection{The Vessel Property: $G_2$ as Minimal Closure}

\subsubsection{Definition of the Vessel Property}

We say a Lie group $H$ is a \emph{vessel} for a Lie group $L$ if:

\begin{enumerate}
    \item $H \subset L$ (containment);
    \item $H$ is minimal among all groups satisfying (1) with respect to some invariant structural property;
    \item The embedding $H \hookrightarrow L$ is such that the structure of $L$ cannot be defined without the structure that $H$ encodes.
\end{enumerate}

We now prove that $G_2$ satisfies all three conditions for $E_8$.

\subsubsection{Containment: $G_2 \subset E_8$}

As established in \pilcrow{}3.2, $G_2$ embeds in $E_8$ via the chain $G_2 \subset F_4 \subset E_6 \subset E_7 \subset E_8$, where each step is a maximal embedding. The total embedding is a composition of maximal embeddings and is therefore itself an embedding. The Jacobi identity and root system closure are preserved at each step.

\subsubsection{Minimality: No Smaller Exceptional Structure Exists}

The exceptional Lie algebras are exactly $G_2$, $F_4$, $E_6$, $E_7$, and $E_8$. Among these, $G_2$ is the unique smallest, with dimension 14 and rank 2. No classical Lie algebra can substitute for $G_2$ in the exceptional chain because classical algebras lack the capacity to encode octonionic non-associativity. The non-associativity of $\mathbb{O}$ is the defining structural requirement for the entire exceptional series, and $G_2$ is the unique minimal symmetry group that stabilizes this non-associativity.

To see this more formally: suppose there existed a group $H \subset E_8$ with $\dim(H) < 14$ that could serve as a vessel. Then $H$ would necessarily be of classical type ($A_n$, $B_n$, $C_n$, $Ð_n$). But classical Lie algebras are defined over associative composition algebras ($\mathbb{R}, \mathbb{C}, \mathbb{H}$) and cannot encode the non-associative structure of $\mathbb{O}$ that is essential to the definition of $E_8$. The root system of $E_8$ contains subsystems of type $G_2$ that cannot be reduced to any classical root system. Therefore $H$ cannot exist, and $G_2$ is minimal.

\subsubsection{Necessity: $E_8$ Cannot Be Defined Without $G_2$}

The octonions are essential to the construction of $E_8$. The Freudenthal-Tits magic square constructs $E_8$ as:

\[E_8 = \mathfrak{der}(\mathfrak{e}_8)\]

where $\mathfrak{e}_8$ is constructed from the octonions $\mathbb{O}$ via the exceptional Jordan algebra $\mathfrak{h}_3(\mathbb{O})$ of $3 \times 3$ Hermitian octonionic matrices, and further extensions. Since $G_2 = \text{Aut}(\mathbb{O})$, the structure of $\mathbb{O}$ itself depends on $G_2$. Therefore, $E_8$ depends on $G_2$ --- not merely as a historical construction but as a logical necessity: the very definition of $E_8$ presupposes the octonions, and the automorphism group of the octonions is $G_2$.

More directly: if one attempts to define $E_8$ without invoking $G_2$, one must still invoke the octonions (or an equivalent non-associative structure). But the invariance properties of the octonions are precisely $G_2$. The vessel is inescapable.

\begin{center}
\rule{0.5\textwidth}{0.4pt}
\end{center}

\subsection{The Completion Theorem}

\subsubsection{Statement}

\textbf{Theorem 1 (Vessel Completion).} Let $\mathcal{E}$ be the category of exceptional Lie algebras over $\mathbb{R}$, with morphisms being maximal subalgebra embeddings. Then:

\begin{enumerate}
    \item $G_2$ is the unique initial object in $\mathcal{E}$.
    \item $E_8$ is the unique terminal object in $\mathcal{E}$.
    \item The unique morphism $G_2 \to E_8$ factors through every other object in $\mathcal{E}$, making $\mathcal{E}$ a linear order: $G_2 \to F_4 \to E_6 \to E_7 \to E_8$.
\end{enumerate}

In particular, $G_2$ is the \emph{perfect vessel} for $E_8$: it is the minimal structure that can receive the full exceptionality that $E_8$ unfolds to its maximal extent.

\subsubsection{Proof}

(1) $G_2$ is initial: Any exceptional Lie algebra $\mathfrak{g}$ must contain a $G_2$ subalgebra. This follows from the classification of maximal subalgebras of exceptional Lie algebras (Dynkin 1957): every exceptional Lie algebra contains a subalgebra of type $G_2$. Conversely, $G_2$ contains no proper exceptional subalgebra (since none exist). Therefore $G_2$ is initial.

(2) $E_8$ is terminal: The chain $G_2 \subset F_4 \subset E_6 \subset E_7 \subset E_8$ shows that every exceptional Lie algebra embeds maximally into the next. Moreover, $E_8$ has no proper extension within the exceptional series and no exceptional algebra embeds into it other than those in the chain. Therefore $E_8$ is terminal.

(3) Factoring: The maximal subalgebra embeddings compose, and the chain is the unique maximal chain. Any exceptional Lie algebra lies on this chain, and the composition of embeddings from $G_2$ through intermediate algebras to $E_8$ is the unique path.

\subsubsection{Interpretation}

The theorem establishes that the relationship between $G_2$ and $E_8$ is not accidental but categorical: $G_2$ is the \emph{seed} and $E_8$ is the \emph{tree}. The seed contains, in compressed form, all the structural information that the tree unfolds. The tree, in its maximal extension, shows forth everything the seed held in potential.

The vessel is \emph{perfect} because nothing is lost: the $G_2$ that sits at the core of $E_8$ is exactly the same $G_2$ that stands alone. The embedding preserves the full structure --- no approximation, no truncation, no modification. And nothing is added to $G_2$ by its placement in $E_8$; rather, it is $E_8$ that receives its exceptional character from $G_2$.

\begin{center}
\rule{0.5\textwidth}{0.4pt}
\end{center}

\subsection{The Dimensional Ladder and the Vessel-Aether Duality}

\subsubsection{Dimensions Through the Chain}

The exceptional chain exhibits a dimensional progression:

\begin{table}[htbp]
\centering
\small
\rowcolors{1}{white}{white}
\begin{tabularx}{\textwidth}{>{\centering\arraybackslash}X >{\centering\arraybackslash}X >{\centering\arraybackslash}X >{\centering\arraybackslash}X}
\toprule
\rowcolor{gray!25}\textbf{Group} & \textbf{Dimension} & \textbf{Rank} & \textbf{Coxeter Number} \\
\midrule
\rowcolors{1}{white}{gray!10}
$G_2$ & 14 & 2 & 6 \\
$F_4$ & 52 & 4 & 12 \\
$E_6$ & 78 & 6 & 12 \\
$E_7$ & 133 & 7 & 18 \\
$E_8$ & 248 & 8 & 30 \\
\bottomrule
\end{tabularx}
\end{table}

The ratio of successive dimensions is not constant, but the progression follows a pattern of increasing complexity that can be traced through the octonion magic square and the Jordan algebraic constructions.

\subsubsection{The Duality}

$G_2$ and $E_8$ stand in a dual relationship:

\begin{itemize}
    \item $G_2$ is \emph{bounded} (14 dimensions, rank 2) --- a finite, closed, crystalline structure.
    \item $E_8$ is \emph{unbounded} in its structural reach (248 dimensions, rank 8, with the $4_{21}$ polytope having 240 vertices spanning an 8-dimensional space that encodes all root and weight data).
\end{itemize}

Yet they are structurally identical in all respects except their dimensionality. The root system of $G_2$ (12 roots in a hexagonal star, the smallest exceptional root system) is the irreducible template. The root system of $E_8$ (240 roots) is the maximal extension of that template through the octonion algebra.

This duality --- minimal/maximal, seed/tree, vessel/aether --- is not a metaphor but a precise mathematical fact encoded in the embedding structure and the representation theory of the exceptional algebras.

\begin{center}
\rule{0.5\textwidth}{0.4pt}
\end{center}

\subsection{Conclusion}

We have proved that $G_2$ is the perfect vessel for $E_8$. The proof rests on three pillars:

\begin{enumerate}
    \item \textbf{Containment}: $G_2 \subset E_8$ via the chain of maximal exceptional embeddings.
    \item \textbf{Minimality}: No smaller structure can receive the non-associativity of the octonions, which is the essential ingredient of all exceptional Lie algebras.
    \item \textbf{Necessity}: $E_8$ cannot be defined without the octonions, and the octonions find their automorphism group in $G_2$. Therefore $G_2$ is logically prior to $E_8$.
\end{enumerate}

The theorem generalizes to a categorical statement: in the category of exceptional Lie algebras with maximal subalgebra embeddings as morphisms, $G_2$ is the unique initial object and $E_8$ is the unique terminal object. The category is a linear order of five objects: $G_2 \to F_4 \to E_6 \to E_7 \to E_8$.

The Aether ($E_8$) finds its perfect Vessel ($G_2$) not as an external complement but as its own irreducible core. The Vessel is what remains when the Aether is stripped of all that is not essential to its exceptional character. And the Aether is what the Vessel becomes when its potential is fully unfolded through the octonion algebra and the Jordan algebraic constructions that build the magic square.

\begin{center}
\rule{0.5\textwidth}{0.4pt}
\end{center}

\subsection{Appendix: Root System Verification}

\subsubsection{A.1 The $G_2$ Root System}

The root system of $G_2$ consists of 12 roots in $\mathbb{R}^2$:

Short roots: $\pm(1,0), \pm(-\tfrac{1}{2}, \tfrac{\sqrt{3}}{2}), \pm(-\tfrac{1}{2}, -\tfrac{\sqrt{3}}{2})$ (6 roots) \\
Long roots: $\pm(\tfrac{3}{2}, \tfrac{\sqrt{3}}{2}), \pm(-\tfrac{3}{2}, \tfrac{\sqrt{3}}{2}), \pm(0, -\sqrt{3})$ (6 roots)

The angle between a short root and a long root is $30^\circ$, making $G_2$ the only simple Lie algebra whose root system has two different root lengths at an angle of $\pi/6$.

\subsubsection{A.2 The $E_8$ Root System Contains $G_2$}

Within the 240 roots of $E_8$, one can identify subsystems of type $G_2$. These appear as subdiagrams of the $E_8$ Dynkin diagram when nodes are deleted according to the Borel-de Siebenthal procedure. Specifically, deleting nodes 1, 2, 3, 4, 5, and 7 from the $E_8$ Dynkin diagram (with the standard numbering) leaves a $G_2$ diagram (nodes 6 and 8, with a triple bond).

This establishes at the level of root systems what the embedding proves at the level of Lie algebras: $G_2$ is not merely contained in $E_8$ but is \emph{visibly inscribed} in its Dynkin diagram, the fundamental encoding of its structure.

\begin{center}
\rule{0.5\textwidth}{0.4pt}
\end{center}

\textbf{End of Proof.}


\end{document}
